\documentclass[11pt]{article}

\usepackage[margin=1in]{geometry}
\usepackage[T1]{fontenc}
\usepackage[utf8]{inputenc}
\usepackage{amsmath,amssymb}
\usepackage{graphicx}
\usepackage{booktabs}
\usepackage[round]{natbib}
\usepackage{hyperref}
\hypersetup{colorlinks=true, linkcolor=blue!50!black, citecolor=blue!50!black, urlcolor=blue!50!black}
\usepackage{enumitem}
\usepackage{xcolor}
\usepackage{caption}
\usepackage{subcaption}

\title{\LARGE Prompt-Driven Interactive Segmentation\\
\large as a Dataset-Annotation Tool\\[4pt]
\large Final Report}
\author{
  \begin{tabular}{c}
    Student A: 2301498\\
    Abdelhafeez Aref Abdelhafeez Almaz
  \end{tabular}
  \and
  \begin{tabular}{c}
    Student B: 2501925\\
    Abdelrahman Eid Ashour Mohamed
  \end{tabular}
}
\date{Milestone 3 --- November 1, 2026\\CSE 681 --- Digital Image Processing}

\begin{document}
\maketitle

\begin{abstract}
Manual instance-mask labeling is labor-intensive, motivating interactive
segmentation tools that use a small number of human prompts. We implemented a
local full-stack annotation application with a React/TypeScript frontend and
FastAPI backend, using a frozen SAM 2.1 Hiera-Tiny checkpoint. The reproducible
data workflow prepares 300 images from COCO 2017 val2017, with 979 non-crowd
instances in 80 categories. The application supports positive and negative
points, boxes, saving accepted masks, and computing IoU and Dice against the
selected COCO instance. A real API smoke test on Apple MPS returned a
$428\times640$ mask with model confidence 0.954; the first request took
1.56 seconds including predictor initialization. However, no manual polygon
baseline was measured and no prompted mask was scored in the available
evaluation manifest. Thus, this work establishes a runnable annotation and
measurement pipeline, but does not establish accuracy, time savings, or
superiority over manual annotation under the tested local setup.
\end{abstract}

\section{Introduction}
\label{sec:intro}
Creating instance masks by hand requires annotators to delineate object
boundaries, which can be costly when preparing image datasets. Interactive
segmentation offers a different workflow: an annotator provides points or a
box, inspects a model-generated mask, and decides whether to refine or accept
it. This project set out to build such a workflow around SAM 2 and evaluate it
against existing COCO instance masks and a manual polygon baseline.

The deliverable is a local web application backed by Python/FastAPI and a
React/TypeScript frontend. It provides a reproducible 300-image COCO 2017
validation subset, target-object selection, point and box prompts, local SAM 2
inference, accepted-mask persistence, and IoU/Dice computation. The software
path was validated with automated tests and one real API inference on Apple
MPS. That inference took 1.56 seconds on the first request, including model
initialization, and returned model confidence 0.954. These are engineering and
smoke-test observations, not comparative evaluation results. The planned
manual polygon experiment and ground-truth-scored prompted sample were not
completed, so the evidence does not support a claim that the tool improves
mask quality or annotation speed.

\section{Related Work}
\label{sec:related}
\textbf{Interactive segmentation.}
GrabCut is a classic foreground-extraction technique in which user-supplied
initialization guides iterative graph-cut refinement \citep{rother2004grabcut}.
This project follows the broader human-guided segmentation idea, but uses a
promptable pretrained model rather than a graph-cut optimization.

\textbf{Promptable foundation models.}
The Segment Anything Model introduced a large-scale promptable segmentation
approach with point and box inputs \citep{kirillov2023sam}. SAM 2 extends this
family to image and video segmentation \citep{ravi2024sam2}; this project uses
its image predictor with the Hiera-Tiny checkpoint and does not train or
fine-tune the model. FastSAM is another recent approach to segment-anything
behavior, using a different model pipeline \citep{zhao2023fastsam}. These works
provide context rather than directly comparable performance numbers: our
smoke-test protocol and sample are not a benchmark reproduction.

\textbf{Dataset and evaluation context.}
COCO supplies images and instance annotations for object recognition and
segmentation research \citep{lin2014coco}. We use its existing instance masks
as references for post-annotation scoring. Each source image remains subject
to its individual license and the dataset terms \citep{cocoTerms}.

\section{Method}
\label{sec:method}
The frozen SAM 2.1 Hiera-Tiny image predictor is wrapped by a FastAPI service.
The React/TypeScript frontend lists local COCO images and their non-crowd
object categories, serves a selected image, and provides positive point,
negative point, and box prompting. A segmentation request sends the image and
prompt coordinates to the local backend. The backend validates the prompt
against image dimensions, runs the image predictor, and returns the
highest-scoring mask among the predictor's multimask candidates, together with
its model confidence. The annotator can inspect the overlay and accept or
refine it.

For an accepted COCO annotation, the backend validates that the target
annotation belongs to the selected image, checks the prediction dimensions,
loads the corresponding exported ground-truth mask, and computes binary IoU
and Dice. It saves the prediction mask, prompt record, per-annotation JSON
record, elapsed time, and an evaluation CSV manifest under the local,
Git-ignored \texttt{data/evaluation/} directory. The timer begins with the
first prompt and ends on acceptance; model inference latency is measured
separately. Ground truth is reserved for scoring and is not shown to guide
prompt selection.

No model training, fine-tuning, or task-specific loss function is used. The
model parameters are those of the downloaded SAM 2.1 checkpoint; prompt
coordinates are supplied directly by the annotator. The high-level pipeline
is:

\begin{figure}[ht]
\centering
\fbox{\parbox{0.94\linewidth}{\centering
COCO image and hidden reference mask
$\longrightarrow$ choose target object
$\longrightarrow$ human point/box prompts
$\longrightarrow$ local SAM 2.1 predictor
$\longrightarrow$ review and accept mask
$\longrightarrow$ persist mask, prompts, and elapsed time
$\longrightarrow$ compare with reference using IoU and Dice}}
\caption{Implemented annotation and evaluation pipeline. The ground-truth
mask is used at the final scoring stage, not as prompt guidance.}
\label{fig:arch}
\end{figure}

\section{Experimental Setup}
\label{sec:setup}

\subsection{Dataset}
The data workflow selects 300 images from COCO 2017 val2017. Eligible images
have between one and eight non-crowd instance annotations; candidates are
sorted before selection with pseudorandom seed 42. The resulting subset
contains 979 non-crowd annotations in 80 categories. Per-instance binary
masks are exported for scoring. The selected image IDs are recorded in
\texttt{docs/selected\_image\_ids.json}; download and regeneration instructions
are in \texttt{data/README.md}. Images, annotation data, masks, and weights are
local-only and ignored by Git. COCO source-image licensing must be checked
before redistribution \citep{cocoTerms}.

\subsection{Baseline(s)}
The planned primary baseline is manual polygon delineation of the same target
instances, using a non-SAM annotation interface and the same timer definition.
The baseline experiment was not run: no manual polygon masks or timing records
were available when this report was prepared. Consequently, this report cannot
confirm a baseline reproduction or compare against published baseline numbers.
The existing COCO masks are reference annotations for scoring, not a measured
manual-tool baseline.

\subsection{Metrics \& Protocol}
For binary prediction mask $P$ and reference mask $G$, intersection over union
and Dice are computed as
\[
\operatorname{IoU}(P,G)=\frac{|P\cap G|}{|P\cup G|},\qquad
\operatorname{Dice}(P,G)=\frac{2|P\cap G|}{|P|+|G|}.
\]
The implementation rejects mismatched mask dimensions and associates each
prediction with its COCO annotation ID. Annotation time is the elapsed
wall-clock time from the first prompt until acceptance; inference latency is
reported separately. A valid paired comparison would require the same
preselected target instances in both workflows, with ground truth withheld
while each mask is created. No paired sample, repeated seed, or statistical
test was completed. Therefore, no aggregate method metrics or time-saving
estimate can be reported.

\section{Results}
\label{sec:results}
The primary experiment remains incomplete. The evaluation directory contained
no accepted predictions or timing manifest when checked for this report, and
no manual-baseline records were available. Accordingly, the two method rows
below explicitly report missing measurements rather than treating missing
data as a zero score.

\begin{table}[ht]
\centering
\caption{Planned paired comparison. Neither workflow has scored observations;
the SAM 2 smoke test below was not accepted or scored as an annotation.}
\label{tab:main}
\begin{tabular}{@{}lccccc@{}}
\toprule
\textbf{Method} & \textbf{Scored $N$} & \textbf{Mean IoU} &
\textbf{Mean Dice} & \textbf{Annotation time} & \textbf{Runtime evidence} \\
\midrule
Manual polygon baseline & 0 & Not measured & Not measured &
Not measured & Not measured \\
SAM 2.1 prompted tool & 0 & Not measured & Not measured &
Not measured & One smoke-test request: 1.56 s \\
\bottomrule
\end{tabular}
\end{table}

The one smoke test used COCO image \texttt{000000000776.jpg} (428 by 640
pixels) and a positive point at $(130,130)$ selected by visual inspection.
With the SAM 2.1 Hiera-Tiny checkpoint on Apple MPS, the first API request
returned a mask with confidence 0.954 in 1.56 seconds. This request included
predictor initialization; it is not steady-state inference latency or
human annotation time. Confidence is not IoU or Dice. No qualitative
success/failure comparison is reported because the mask was not part of a
reviewed, scored annotation trial.

\begin{figure}[ht]
\centering
\fbox{\parbox{0.90\linewidth}{\centering
No paired qualitative examples are reported. The available single smoke-test
mask was not reviewed and scored against its target before this report.}}
\caption{Qualitative comparison unavailable: no verified paired manual and
SAM 2 annotations were collected.}
\label{fig:qual}
\end{figure}

\subsection{Ablation Studies}
\label{sec:ablation}
The planned prompt-type ablation would compare point prompting with box
prompting on the same target instances, holding the checkpoint and acceptance
protocol fixed. It was not performed, so no effect of prompt type can be
estimated.

\begin{table}[ht]
\centering
\caption{Prompt-type ablation status. Missing values indicate that the
experiment was not run.}
\label{tab:ablation}
\begin{tabular}{@{}lccc@{}}
\toprule
\textbf{Prompt condition} & \textbf{Scored $N$} & \textbf{Mean IoU} &
\textbf{Mean annotation time} \\
\midrule
Positive point & 0 & Not measured & Not measured \\
Box & 0 & Not measured & Not measured \\
\bottomrule
\end{tabular}
\end{table}

\section{Discussion \& Limitations}
\label{sec:discussion}
The completed contribution is a reproducible software and data path: the
application can select COCO targets, accept interactive prompts, request local
SAM 2 inference, and persist accepted masks with timing and mask metrics. The
backend suite passed 17 tests, the frontend TypeScript production build and
lint completed successfully, and a real one-image inference request verified
that the downloaded checkpoint can run through the API on Apple MPS. These
checks establish implementation behavior, not segmentation quality.

The central limitation is the absence of the planned paired human study. The
manual polygon baseline was not collected, and no accepted SAM 2 masks were
recorded in the evaluation manifest. The single smoke test has no ground-truth
score and cannot support an accuracy or time-saving claim. In addition, the
sample was designed as a small reproducibility subset rather than a
representative benchmark; no annotator study, inter-annotator analysis,
uncertainty estimate, or per-category analysis was performed. The observed
first-request latency includes model initialization and comes from one
machine/device, so it should not be generalized to other hardware or
steady-state requests.

The system also depends on local model weights and COCO assets that are
excluded from version control. Reproduction requires downloading them with the
documented scripts, installing the optional SAM 2 dependency, and matching the
checkpoint to its configuration. The frontend and backend automate metric
calculation but do not by themselves ensure that annotation prompts were
chosen without seeing ground truth; this remains a protocol requirement for
future trials. The results are therefore an honest implementation milestone,
not evidence that SAM 2 is faster or more accurate than manual polygon
labeling.

\section{Conclusion}
\label{sec:conclusion}
We built a local React/FastAPI annotation tool around frozen SAM 2.1 and a
reproducible 300-image COCO subset, including an API path that stores masks and
computes IoU, Dice, and annotation time. One real inference smoke test and
automated software checks succeeded, but the paired manual baseline and scored
prompted evaluation were not completed. No conclusion about quality or time
savings is therefore justified; the next step is to collect blinded,
same-instance annotations with the shared timer protocol.

\section*{Individual Contributions}
\label{sec:contributions}
Student A was assigned COCO data preparation and the manual polygon baseline;
the baseline measurements remain incomplete. Student B implemented the local
React/FastAPI annotation workflow, SAM 2 integration, metric persistence, and
evaluation tooling.

\bibliographystyle{unsrtnat}
\bibliography{references}

\appendix
\section{Additional Results}
\label{sec:appendix}
The reproducible software checks at report time were 17 passing backend tests,
successful Ruff checks, a successful frontend TypeScript/Vite production
build, and a successful frontend lint run. The real inference smoke test
returned a 428 by 640 mask with confidence 0.954; the first API request took
1.56 seconds on Apple MPS including predictor initialization. These are
validation observations only and are not included as quality metrics.

\end{document}
